\section{Motivation}
\label{sec:whyfinetune}

The privacy goal fixes the shape of the protocol before any modeling choice is made. The contributions must
be combined while encrypted, in a single round, and the party that combines them must learn nothing in the
process. That party is therefore \emph{blind}: it holds ciphertexts that the security of the scheme makes
indistinguishable from random, and it must produce the aggregate without ever seeing a value it operates on.

Blindness is a stronger constraint than cost. A homomorphic server evaluates a fixed arithmetic circuit chosen
in advance, and it cannot read a value, compare two values, or branch on the outcome of a comparison without
the deep circuits, bootstrapping, or extra interaction that a one-shot, low-depth protocol rules out. Little
programmability is available in practice: the server cannot cheaply decide which client to trust, which
coordinates to keep, or how to precondition an update by its magnitude, because each of those decisions would
have to read the encrypted data. The computation a blind server can afford is a shallow, data-independent one,
and the natural such computation is a weighted sum. Whatever adaptivity the method needs must therefore live
where the server is not: in the clients' local plaintext training, or in public scalars --- the sample
weights and the scaling coefficient --- that reveal nothing about any client's data.

A blind weighted sum is only meaningful when the things being summed already agree. Averaging two networks
that began from different initializations and fit different data gives a model worse than either, because
their minima are not joined by a straight line and the combination lands on the ridge between two basins
rather than inside one. For the server's single operation to return a usable model, the clients' models must
share a loss basin: a common point they all departed from by a bounded amount, so the straight line between
them stays inside one minimum. We see this directly in the experiments --- remove the shared basin and the
blind linear combination collapses toward chance, while restoring it makes the same operation yield a strong
model (\cref{sec:necessity}).

Training from scratch cannot supply that common point in a single round. Independently initialized clients
have no shared frame, and the methods that recover one after the fact, by matching neurons or permutations
across models, must read the models --- which is exactly what encryption withholds. A pretrained model
supplies this frame already. Freezing the backbone fixes the representation every client sees, so each
client's trainable parameters start at the same place and move only a short distance; the shared loss basin
\emph{is} the pretrained model. Fine-tuning a small set of parameters on top of a fixed, common backbone is
therefore not a deployment convenience but the learning paradigm that fits a one-shot, blind, low-depth
protocol, and it is what lets the encrypted object be a handful of adapter parameters rather than a whole
network.

This also fixes the rest of the design. The trainable unit is a low-rank adapter and a head rather than the
full model, because keeping the fine-tuned parameters few keeps the circuit shallow and the ciphertexts cheap
while leaving the learning where it belongs, in the client's plaintext fine-tuning. The server's role
contracts to the one thing it can soundly do --- add the clients' bounded fine-tuning displacements, scaled by
public weights --- and that operation is a depth-one homomorphic sum, the task arithmetic of
\cref{sec:method}. The accuracy the frozen features do not already supply is then provided not by the server
but by the clients, whose local fine-tuning is the source of the lift that the aggregate inherits
(\cref{sec:basin-contribution}).
